%! AP Statistics — Unit 2, Lesson 2.2 — Warm-Up KEY
\documentclass[10pt]{article}
\usepackage{apstats-article}
\usepackage{apstats-key}

%=============================================================================
\begin{document}

\pageheader{Unit 2, Lesson 2.2}{Warm-Up --- Answer Key}
\begin{tabularx}{\linewidth}{X X X}
Name:\;\ans{Key} & Date:\; & Period:\;
\end{tabularx}

\vspace{0.3in}

{\large\bfseries\color{navy} Part 1 \quad Frequency Tables (Lesson 1.4 Review)}

\vspace{0.12in}

\noindent A survey asked 60 students which after-school activity they prefer.
The results are shown below.

\vspace{0.1in}

\renewcommand{\arraystretch}{1.6}
\begin{tabularx}{\linewidth}{@{} l C{2.8cm} C{2.8cm} @{}}
\rowcolor{sky}
\textbf{Activity}   & \textbf{Frequency} & \textbf{Relative Frequency} \\
\hline
Sports              & 24 & \ans{$24/60 = 0.40$} \\
Arts / Music        & 18 & \ans{$18/60 = 0.30$} \\
Academic clubs      & 12 & \ans{$12/60 = 0.20$} \\
None / Other        &  6 & \ans{$6/60 = 0.10$} \\
\hline
\textbf{Total}      & \ans{60} & \ans{1.00} \\
\end{tabularx}

\vspace{0.18in}

\begin{enumerate}

    \item \ans{Relative frequency $= 24 \div 60 = 0.40$ (or 40\%).}

    \vspace{0.1in}

    \item \ans{40\% of surveyed students preferred sports.}

    \vspace{0.1in}

    \item \ans{The bar height would be 0.30 (or 30\%), representing the proportion of
    students who preferred Arts/Music out of all 60 surveyed.}

\end{enumerate}

\vspace{0.3in}

{\large\bfseries\color{navy} Part 2 \quad Two Variables at Once}

\vspace{0.12in}

\noindent Now suppose we also recorded whether each student is an underclassman
(9th/10th grade) or upperclassman (11th/12th grade).

\vspace{0.14in}

\begin{tcolorbox}[colback=greenbg, colframe=greenacc, boxrule=0.8pt, arc=2mm,
    left=4mm, right=4mm, top=3mm, bottom=3mm]
    \small
    \textbf{Think about it:} The one-way table above describes one variable at a time.
    If we want to know \emph{whether grade level is related to activity preference},
    we need to cross \textbf{two} categorical variables simultaneously.
    \begin{itemize}[leftmargin=1.3em, itemsep=6pt]
        \item A \textbf{two-way table} organizes two categorical variables at once
              --- you will build one in today's guided notes.
        \item What would the rows represent? What would the columns represent?
    \end{itemize}
\end{tcolorbox}

\vspace{0.16in}

\noindent\textbf{Sketch a blank two-way table.}

\vspace{0.15in}

\renewcommand{\arraystretch}{2.2}
\begin{tabularx}{\linewidth}{@{} l Y Y Y Y Y @{}}
\rowcolor{sky}
\textbf{Grade} & \textbf{Sports} & \textbf{Arts} & \textbf{Clubs} & \textbf{None} & \textbf{Total} \\
\hline
Under          & \ans{} & \ans{} & \ans{} & \ans{} & \ans{} \\
Upper          & \ans{} & \ans{} & \ans{} & \ans{} & \ans{} \\
\hline
Total          & \ans{} & \ans{} & \ans{} & \ans{} & \ans{} \\
\end{tabularx}

\begin{teachernote}
Part 1: Students should compute all four relative frequencies and verify they sum to 1.00.
Common error: dividing by a row total rather than the grand total of 60.
Part 2: Accept any blank two-way table structure with grade level as rows and activity as columns.
The purpose is to activate schema for what a two-way table looks like before guided notes.
\end{teachernote}

\end{document}
